\problemname{Noughts and Crosses}
Simon and Måns are sitting on the subway, playing a game.
Simon draws an $n \times m$ grid and fills in a few crosses and noughts in the cells.
He also writes a number on each row.

Måns's task is now to fill in the grid so that \emph{every} cell has a cross or a nought, such that
there are exactly as many crosses on each row as the number written there.
The goal for Måns is that every row, column, and diagonal has at most $2$ crosses or noughts in a row.

Simon has placed the crosses and noughts completely at random, so it is not obvious that this is possible.
However, we can give Måns points depending on how well he succeeded: we define the value of a placement
as the number of times there are three crosses or three noughts in a row in any row, column, or diagonal
(in both directions).
Måns should then try to get as low a value as he can.

Help Måns play the game as well as possible!

\section*{Input}
\textbf{Note:} the test data for this problem is \emph{open}. It is available as an attachment to the problem.
The first line contains an integer $0 \le t \le 10$, the index of this test case.
The case $t = 0$ represents the sample case and should be ignored (you can output anything then).

The second line contains two integers $n$ and $m$ ($2 \le n, m \le 500$): the height and width of the grid.
The third line contains $n$ integers: the number of crosses that must be placed on each row.
Then follow $n$ lines, with $m$ characters each: the original grid.
Each character will be either ``\texttt{.}'', ``\texttt{o}'' or ``$\texttt{x}$'',
where ``\texttt{.}'' means the cell has not yet been filled in.

It is guaranteed that the number of crosses already in each row is at most equal to the number Simon wrote for that row.

\section*{Output}
Output $n$ lines with $m$ characters each: the fully filled-in grid, where all ``\texttt{.}'' have been replaced
by either ``\texttt{o}'' or ``\texttt{x}''.

Note that Kattis has a source code size limit of 128kB. For grids of maximum size, it is therefore
\emph{not} possible to hardcode complete solutions in the code.

\section*{Scoring}
Your solution will receive points on each of the ten test cases.

Each test case is scored as follows.
Assume that the value you achieved on the test case is $X$, and
the minimum value any contestant achieved during the contest is $X_{\min}$.
If $X_{\min} \ne 0$, let $X_{\min}' = X_{\min}$, otherwise let $X_{\min}'$ be the second best contestant's value on that test case (which is not 0).
If $X \ge 2X_{\min}'$, the test case is awarded $0$ points, otherwise it is awarded $10(1 - r)$ points, where $r = (X - X_{min}) / (2X_{\min}' - X_{\min})$.

During the contest, $X_{min}$ will be set to a value computed by the judges.
This means that the scoreboard may temporarily show points up to 200, and it is possible that your points on a test case may be recalculated to 0 after the contest!

\section*{Explanation of Sample 1}
In the solution to the sample there are 26 vertical three-in-a-rows, 18 horizontal, and 32 diagonal.
In total, this gives a value of 76.
